\documentclass{article}
\usepackage{amsmath,amssymb,amsthm,mathtools}
\usepackage{proofread}
\newtheorem{theorem}{Theorem}
\newtheorem{lemma}[theorem]{Lemma}
\theoremstyle{remark}
\newtheorem{remark}[theorem]{Remark}

\title{The harmonic, geometric and arithmetic mean of two numbers}
\author{Proofread example}
\date{}

\begin{document}
\maketitle

\begin{abstract}
For two positive numbers, the harmonic mean is at most the geometric mean, and the geometric mean is at most the arithmetic mean.
The geometric and the arithmetic mean agree exactly when the two numbers do.
We give the short classical proof, in which every step is derived from what it cites.
\end{abstract}

\section{Introduction}

The inequality between the arithmetic and the geometric mean is one of the most frequently used inequalities in analysis, and the monographs \cite{hardy1952inequalities,steele2004cauchy} discuss it at length.
For two numbers it reduces to the fact that a square is nonnegative.
The inequality for the harmonic mean then follows from a short computation.
Theorem~\ref{thm:means} states both, together with the case of equality.

\section{The three means}

\begin{lemma}\label{lem:square}
Let $a,b\in(0,\infty)$.
Then $\frac{a+b}{2}-\sqrt{ab}=\frac{d^2}{2}$, where $d\coloneqq\sqrt a-\sqrt b$.
\end{lemma}

\begin{proof}
We have $(\sqrt a)^2=a$, $(\sqrt b)^2=b$ and $\sqrt a\sqrt b=\sqrt{ab}$.
Hence $\frac{d^2}{2}=\frac{a-2\sqrt{ab}+b}{2}=\frac{a+b}{2}-\sqrt{ab}$.
\end{proof}

\begin{theorem}\label{thm:means}
Let $a,b\in(0,\infty)$, and let $A\coloneqq\frac{a+b}{2}$, $G\coloneqq\sqrt{ab}$ and $H\coloneqq\frac{2ab}{a+b}$.
Then $H\le G\le A$.
Moreover, $G=A$ if and only if $a=b$.
\end{theorem}

\begin{proof}
By Lemma~\ref{lem:square}, $A-G=\frac{(\sqrt a-\sqrt b)^2}{2}\ge0$.
Hence $G\le A$, and $G=A$ if and only if $\sqrt a=\sqrt b$.
Moreover, $\sqrt a=\sqrt b$ if and only if $a=b$, since the square root is injective.
Finally, $H=\frac{G^2}{A}$.
Since $G\le A$ and $G>0$, it follows that $H=G\cdot\frac{G}{A}\le G$.
\end{proof}

\begin{remark}\label{rem:quadratic}
For $a,b\in(0,\infty)$, the quadratic mean $Q\coloneqq\sqrt{\frac{a^2+b^2}{2}}$ is at least the arithmetic mean $\frac{a+b}{2}$.
\end{remark}

\begin{thebibliography}{2}
\bibitem{hardy1952inequalities} G.~H. Hardy, J.~E. Littlewood and G.~P\'olya, \emph{Inequalities}, 2nd ed., Cambridge University Press, 1952.
\bibitem{steele2004cauchy} J.~M. Steele, \emph{The Cauchy--Schwarz Master Class}, Cambridge University Press, 2004.
\end{thebibliography}

\end{document}
